\scanpage{17}
Let us likewise eliminate the case $g=0$, $\nu\leq2$. Here $g=0=\sum g_\alpha+h_1$ implies that the $g_\alpha$ and $h_1$ are zero : all the vertices are of genus $0$, the graph is a tree. It cannot be reduced to a single vertex, so it has at least two ends, whose marked \unsure{weights|points} must be $\geq2$, so \deleted{\unsure{the}} we have $\nu\geq2+2=4$, absurd.

Let $G$ be a \deleted{\inserted{MD\unsure{-}}graph}\inserted{maquette}. One says that a \deleted{\inserted{MD\unsure{-}}} \deleted{\unsure{standart|stable}}\inserted{standart} curve $C$ over an algebraically closed field is of \uline{type $G$}, if its maquette is isomorphic to $G$, one says that it is \inserted{$G$-}\uline{pinned} if one gives oneself an isomorphism between its maquette and $G$ (this is therefore an additional\unsure{-|.} structure on $\unsure{X|C}$, called a $G$-pinning).

Let $(\widehat X,\underline I)$ be a standart curve over an arbitrary base $\mathcal S$, one says that it is of type $G$ if its geometric fibers are of type $G$. Then the \deleted{\unsure{``types'' of the}}\inserted{maquettes} of the
% Light lower-left pagination; hand unidentified.
\mspage{16}
